\chapter{路易斯结构、共振与价层电子对互斥模型}\label{ch:b1:lewis-resonance-vsepr}

臭氧分子 \ce{O3} 是三个氧原子排成的折线形链。
若按第 1 册的规则画出它，会得到一个双键和一个单键，
于是一根键应较短，另一根应较长。然而微波光谱测得，
臭氧中的两根 \ce{O-O} 键完全相等，键长为 \qty{127.8}{pm}，
介于 \ce{O2} 的双键（\qty{120.8}{pm}）与过氧化氢的单键（\qty{147.5}{pm}）之间。
没有任何一种画法是正确的，两种画法合在一起才正确。
本章使路易斯模型更加精确——包括\omterm{def:b1:lewis-resonance-vsepr:formal-charge}{形式电荷}、
不完整或超出的八隅体，以及若干结构之间的共振——
然后用它预测分子的形状，
以及由此而来的\omterm{def:b1:lewis-resonance-vsepr:dipole}{偶极矩}。
% ledger: geo:O3.rOO, re:O2, geo:H2O2.rOO

\begin{recall}
第 1 册（十年级）描述了\omterm{def:b1:lewis-resonance-vsepr:lewis}{共价键}是一对共用电子，
画出了带有\omterm{def:b1:lewis-resonance-vsepr:lewis}{成键电子对}和\omterm{def:b1:lewis-resonance-vsepr:lewis}{孤电子对}的\omterm{def:b1:lewis-resonance-vsepr:lewis}{路易斯结构}，
并预测了四种形状（直线形、折线形、三角形、四面体形）。
\omterm{def:b1:periodicity:electronegativity}{电负性}及其差值已在\cref{ch:b1:periodicity}中作了定量化；
原子的\omterm{def:b1:quantum-numbers:configuration}{价电子}可由其\omterm{def:b1:quantum-numbers:configuration}{电子排布}读出
（见\cref{def:b1:quantum-numbers:configuration}）。
\end{recall}

\section{路易斯结构与形式电荷}

\begin{definition}[共价键，路易斯结构]\label{def:b1:lewis-resonance-vsepr:lewis}
两个原子共用的一对电子称为
\emph{共价键}\index{共价键}，即\emph{成键电子对}\index{成键电子对}；
两对或三对成键电子对分别构成双键或三键。
属于单个原子的一对\omterm{def:b1:quantum-numbers:configuration}{价电子}称为\emph{孤电子对}\index{孤电子对}。
\emph{路易斯结构}\index{路易斯结构}表示分子或离子的全部\omterm{def:b1:quantum-numbers:configuration}{价电子}，
其中成键电子以原子之间的连线表示，
孤电子对以原子上的线或点对表示。
根据\emph{八隅规则}\index{八隅规则}，第二\omterm{def:b1:periodicity:table}{周期}的原子（C、N、O、F）
在稳定结构中被四对即八个电子包围；
根据\emph{二隅规则}\index{二隅规则}，氢原子被一对电子包围。
\end{definition}

\begin{definition}[形式电荷]\label{def:b1:lewis-resonance-vsepr:formal-charge}
在\omterm{def:b1:lewis-resonance-vsepr:lewis}{路易斯结构}中，某原子的\emph{形式电荷}\index{形式电荷}
为
\[
  q_F = N_v - N_{\text{孤对}} - \tfrac12 N_{\text{成键}},
\]
其中 $N_v$ 为游离原子的\omterm{def:b1:quantum-numbers:configuration}{价电子}数，
$N_{\text{孤对}}$ 为其\omterm{def:b1:lewis-resonance-vsepr:lewis}{孤电子对}中的电子数，
$N_{\text{成键}}$ 为它所形成的键中的电子数。
当其不为零时，以 $\oplus$ 或 $\ominus$ 写在该原子旁边。
\end{definition}

\begin{proposition}[形式电荷之和等于总电荷]\label{prop:b1:lewis-resonance-vsepr:formal-sum}
\omterm{def:b1:lewis-resonance-vsepr:lewis}{路易斯结构}中各原子的\omterm{def:b1:lewis-resonance-vsepr:formal-charge}{形式电荷}之和，
等于分子或离子的总电荷。
\end{proposition}

\begin{proof}
对各原子的 $q_F$ 求和，得 $\sum N_v - (\text{孤电子对中的
电子数} + \text{键中的电子数})$，这是因为每根键的电子
在其两个原子上各计一半。括号中的量是结构中画出的
\omterm{def:b1:quantum-numbers:configuration}{价电子}总数，它等于 $\sum N_v$ 减去该物种的电荷 $z$，
因此 $\sum q_F = z$。
\end{proof}

\begin{method}[画路易斯结构]\label{met:b1:lewis-resonance-vsepr:lewis}
画出分子或离子的\omterm{def:b1:lewis-resonance-vsepr:lewis}{路易斯结构}，步骤如下：
\begin{enumerate}
  \item 计算\omterm{def:b1:quantum-numbers:configuration}{价电子}数 $N = \sum N_v - z$ 及电子对数 $N/2$；
  \item 画出骨架：\omterm{def:b1:periodicity:electronegativity}{电负性}最小的原子（氢除外）通常居中；
        氢和氟总是处于端位；
  \item 用\omterm{def:b1:lewis-resonance-vsepr:lewis}{孤电子对}补全端位原子的八隅体，然后
        将剩余电子对放在中心原子上；
  \item 若中心原子缺少八隅体，则将其相邻原子的\omterm{def:b1:lewis-resonance-vsepr:lewis}{孤电子对}
        转化为多重键；
  \item 计算\omterm{def:b1:lewis-resonance-vsepr:formal-charge}{形式电荷}；在可能的结构中，优先选择
        \omterm{def:b1:lewis-resonance-vsepr:formal-charge}{形式电荷}最少、且负电荷位于
        \omterm{def:b1:periodicity:electronegativity}{电负性}最大的原子上的结构。
\end{enumerate}
\end{method}

\begin{example}[硝酸]\label{ex:b1:lewis-resonance-vsepr:hno3}
\ce{HNO3}：$N = 1 + 5 + 3 \times 6 = 24$ 个电子，即 12 对。骨架：
氮居中，与三个氧原子相连，其中一个氧原子连有氢。
若以单键补全八隅体，则氮只有六个电子；
将一个氧的\omterm{def:b1:lewis-resonance-vsepr:lewis}{孤电子对}转化为 \ce{N=O} 键。\omterm{def:b1:lewis-resonance-vsepr:formal-charge}{形式电荷}：氮为
$5 - 0 - 4 = +1$；不连氢的单键氧为
$6 - 6 - 1 = -1$；其余为 0。总电荷为 0，符合
\cref{prop:b1:lewis-resonance-vsepr:formal-sum}的要求。
\end{example}

\begin{omfigure}
\setchemfig{atom sep=2.4em}
\begin{tikzpicture}[font=\small]
  \node at (0,0) {\large\chemfig{H-\charge{90=\|,270=\|}{O}-\charge{90:2pt=$\scriptstyle\oplus$}{N}(=[:45]\charge{0=\|,90=\|}{O})-[:-45]\charge{0=\|,240=\|,300=\|,45:3pt=$\scriptstyle\ominus$}{O}}};
  \node[below] at (0,-1.5) {硝酸，\ce{HNO3}};
  \node at (5.2,0) {\large\chemfig{\charge{180=\|}{N}~\charge{0=\|}{N}}};
  \node[below] at (5.2,-1.5) {氮气，\ce{N2}};
  \node at (9.6,0) {\large\chemfig{\charge{135=\|,225=\|,90:2pt=$\scriptstyle\ominus$}{N}=\charge{90:2pt=$\scriptstyle\oplus$}{N}=\charge{45=\|,315=\|}{O}}};
  \node[below] at (9.6,-1.5) {一氧化二氮，\ce{N2O}};
\end{tikzpicture}

{\small 带\omterm{def:b1:lewis-resonance-vsepr:lewis}{孤电子对}与\omterm{def:b1:lewis-resonance-vsepr:formal-charge}{形式电荷}的\omterm{def:b1:lewis-resonance-vsepr:lewis}{路易斯结构}。}
\end{omfigure}

\section{超出八隅体}

\begin{definition}[缺电子分子与超价分子]\label{def:b1:lewis-resonance-vsepr:hypervalent}
若分子中某原子的\omterm{def:b1:quantum-numbers:configuration}{价电子}少于八个，则称该分子为
\emph{缺电子的}\index{缺电子的}分子：在 \ce{BF3} 中，硼有六个\omterm{def:b1:quantum-numbers:configuration}{价电子}。
若第三\omterm{def:b1:periodicity:table}{周期}及以后的原子周围的电子对多于四对，
则称该分子为\emph{超价的}\index{超价的}分子：在 \ce{PCl5} 中，磷形成五根键，
在 \ce{SF6} 中，硫形成六根键。含奇数个电子的分子，例如 \ce{NO} 和 \ce{NO2}，
不能使每个原子都满足\omterm{def:b1:lewis-resonance-vsepr:lewis}{八隅规则}，
其中有一个电子未成对。
\end{definition}

\begin{remark}[为何是第三周期而非第二周期]\label{rem:b1:lewis-resonance-vsepr:hypervalent}
氮能形成 \ce{NF3}，却不能形成 \ce{NF5}，而磷则 \ce{PF3} 和 \ce{PF5} 都能形成。
第二\omterm{def:b1:periodicity:table}{周期}原子太小，不能与四个以上的邻原子成键；
第三\omterm{def:b1:periodicity:table}{周期}原子足够大，能容纳五个或六个。
\ce{SO4^{2-}} 或 \ce{H2SO4} 的结构式，既有画成硫与氧之间为双键（硫周围有 12 个电子）的，
也有画成单键并标出\omterm{def:b1:lewis-resonance-vsepr:formal-charge}{形式电荷}（八隅体）的，两者都能见到；
后者更好地反映了电荷分布，
两种画法都给出正确的几何构型。
\end{remark}

\begin{definition}[路易斯酸、路易斯碱与配位键]\label{def:b1:lewis-resonance-vsepr:lewis-acid}
称\emph{路易斯酸}\index{路易斯酸}为能够将一对电子接收到其价层空位中的物种
（如 \ce{BF3}、\ce{AlCl3}、\ce{H+}、金属阳离子）；
带有可共享的\omterm{def:b1:lewis-resonance-vsepr:lewis}{孤电子对}的物种，称为\emph{路易斯碱}\index{路易斯碱}
（如 \ce{NH3}、\ce{H2O}、
\ce{F-}）。当碱将电子对给予酸时，所形成的键称为
\emph{配位键}\index{配位键}：这是一种\omterm{def:b1:lewis-resonance-vsepr:lewis}{共价键}，
其两个电子都来自同一方。
\end{definition}

\begin{example}[氨与三氟化硼]\label{ex:b1:lewis-resonance-vsepr:adduct}
\ce{NH3 + BF3 -> H3NBF3}。加合物中，硼达到八隅体；
氮现已提供其\omterm{def:b1:lewis-resonance-vsepr:lewis}{孤电子对}，带\omterm{def:b1:lewis-resonance-vsepr:formal-charge}{形式电荷} $+1$，
硼带\omterm{def:b1:lewis-resonance-vsepr:formal-charge}{形式电荷} $-1$。一旦形成，\ce{N-B} 键
便是普通的\omterm{def:b1:lewis-resonance-vsepr:lewis}{共价键}。
\end{example}

\begin{omfigure}
\setchemfig{atom sep=2.1em}
\begin{tikzpicture}[font=\small]
  \node at (0,0) {\large\chemfig{H-\charge{90=\|}{N}(-[:240]H)(-[:300]H)}};
  \node at (1.6,0) {$+$};
  \node at (3.1,0) {\large\chemfig{B(-[:90]F)(-[:210]F)(-[:330]F)}};
  \draw[-{Stealth}, thick] (4.3,0) -- (5.3,0);
  \node at (7.4,0) {\large\chemfig{H-\charge{45:1pt=$\scriptstyle\oplus$}{N}(-[:90]H)(-[:270]H)-\charge{45:1pt=$\scriptstyle\ominus$}{B}(-[:90]F)(-[:270]F)-F}};
\end{tikzpicture}

{\small \omterm{def:b1:lewis-resonance-vsepr:lewis-acid}{配位键}：氨的\omterm{def:b1:lewis-resonance-vsepr:lewis}{孤电子对}（\omterm{def:b1:lewis-resonance-vsepr:lewis-acid}{路易斯碱}）
填入三氟化硼的空位（\omterm{def:b1:lewis-resonance-vsepr:lewis-acid}{路易斯酸}）。
氟的\omterm{def:b1:lewis-resonance-vsepr:lewis}{孤电子对}未画出。}
\end{omfigure}

\section{共振}

\begin{definition}[共振结构与共振杂化体]\label{def:b1:lewis-resonance-vsepr:resonance}
若分子的电子可以放入若干个\omterm{def:b1:lewis-resonance-vsepr:lewis}{路易斯结构}中，这些结构仅在 $\pi$ 键和\omterm{def:b1:lewis-resonance-vsepr:lewis}{孤电子对}的位置上不同
——原子核保持不动——则每一个都称为一个\emph{共振结构}\index{共振结构}
（或称介变形式），它们之间以双箭头 $\leftrightarrow$ 相连。
真实分子并非其中任何一个，
而是\emph{共振杂化体}\index{共振杂化体}，
是一个单一的结构，
其电子分布为各共振结构的
加权平均。
\end{definition}

\begin{remark}[是箭头，而非平衡]\label{rem:b1:lewis-resonance-vsepr:arrow}
箭头 $\leftrightarrow$ 并不表示反应：臭氧并不会
从一个结构转变为另一个结构。\omterm{def:b1:lewis-resonance-vsepr:hybridisation}{杂化}体是同一个分子，
之所以用若干图形来描述，是因为一个图形不足以说明。
\omterm{def:b1:lewis-resonance-vsepr:resonance}{共振结构}之间决不可使用 $\rightleftharpoons$。
\end{remark}

\begin{method}[书写与加权共振结构]\label{met:b1:lewis-resonance-vsepr:resonance}
求某一物种的\omterm{def:b1:lewis-resonance-vsepr:resonance}{共振结构}，步骤如下：
\begin{enumerate}
  \item 从一个\omterm{def:b1:lewis-resonance-vsepr:lewis}{路易斯结构}出发；寻找\omterm{def:b1:lewis-resonance-vsepr:lewis}{孤电子对}、
        紧邻另一 $\pi$ 键的 $\pi$ 键、
        空位或正电荷；
  \item 用弯箭头移动电子对（\omterm{def:b1:lewis-resonance-vsepr:lewis}{孤电子对}变为 $\pi$ 键，
        $\pi$ 键变为\omterm{def:b1:lewis-resonance-vsepr:lewis}{孤电子对}），决不移动原子核，
        也决不使第二\omterm{def:b1:periodicity:table}{周期}原子超过八隅体；
  \item 重新计算\omterm{def:b1:lewis-resonance-vsepr:formal-charge}{形式电荷}；
  \item 为各结构加权：最重要的结构满足\omterm{def:b1:lewis-resonance-vsepr:lewis}{八隅规则}，
        \omterm{def:b1:lewis-resonance-vsepr:formal-charge}{形式电荷}最少，
        并将负电荷置于\omterm{def:b1:periodicity:electronegativity}{电负性}最大的原子上；
        等价结构权重相同。
\end{enumerate}
\end{method}

\begin{example}[臭氧、硝酸根与苯]\label{ex:b1:lewis-resonance-vsepr:examples}
臭氧有两个等价结构：每个 \ce{O-O} 键在一个结构中为双键，
在另一个结构中为单键，因此两者的键级均为 $1.5$，
长度相同，均为 \qty{127.8}{pm}，介于 \ce{O=O} 与 \ce{O-O} 之间。
硝酸根离子 \ce{NO3-} 有三个等价结构；每个 \ce{N-O} 键的键级为 $4/3$。
苯 \ce{C6H6} 有两个凯库勒结构；其六个 \ce{C-C} 键相等，
键长均为 \qty{139.7}{pm}，介于乙烷的单键
（\qty{153.6}{pm}）与乙烯的双键（\qty{133.9}{pm}）之间。
% ledger: geo:O3.rOO, geo:C6H6.rCC, geo:C2H6.rCC, geo:C2H4.rCC
\end{example}

\begin{omfigure}
\vspace*{10pt}
\large
\setchemfig{atom sep=2.4em}
\schemestart
\chemfig{\charge{135=\|,225=\|}{O}=[:-30]\charge{270=\|,90:2pt=$\scriptstyle\oplus$}{O}-[:30]\charge{90=\|,0=\|,270=\|,45:5pt=$\scriptstyle\ominus$}{O}}
\arrow{<->}
\chemfig{\charge{90=\|,180=\|,270=\|,135:5pt=$\scriptstyle\ominus$}{O}-[:-30]\charge{270=\|,90:2pt=$\scriptstyle\oplus$}{O}=[:30]\charge{45=\|,315=\|}{O}}
\schemestop

\medskip
\schemestart
\chemfig{*6(=-=-=-)}
\arrow{<->}
\chemfig{*6(-=-=-=)}
\schemestop
\hspace{3em}
\chemfig{\charge{270:2pt=$\scriptstyle\oplus$}{N}(=[:90]\charge{45=\|,135=\|}{O})(-[:210]\charge{150=\|,240=\|,300=\|,90:2pt=$\scriptstyle\ominus$}{O})-[:330]\charge{30=\|,300=\|,240=\|,90:2pt=$\scriptstyle\ominus$}{O}}
\normalsize
\par\vspace{14pt}
{\small 臭氧中的共振（右侧氧的一对\omterm{def:b1:lewis-resonance-vsepr:lewis}{孤电子对}转化为 $\pi$ 键，
同时左侧的 $\pi$ 键转化为\omterm{def:b1:lewis-resonance-vsepr:lewis}{孤电子对}）与苯中的共振，
以及硝酸根离子三个等价结构之一，
其电荷
由三个氧原子共同分担。}
\end{omfigure}

\begin{definition}[$\sigma$ 键、$\pi$ 键]\label{def:b1:lewis-resonance-vsepr:sigma-pi}
单键称为\emph{$\sigma$ 键}\index{$\sigma$ 键}：
其电子对沿连接两个原子核的轴线分布，
由两个轨道头对头重叠形成。多重键中的第二根和第三根键
称为\emph{$\pi$ 键}\index{$\pi$ 键}：其电子对位于轴线两侧，
由两个 p 轨道肩并肩重叠形成，这两个轨道
均与轴线垂直。双键由一根 $\sigma$ 键和一根 $\pi$ 键构成；
三键由一根 $\sigma$ 键和两根 $\pi$ 键构成。
\end{definition}

\begin{omfigure}
\begin{tikzpicture}[font=\small]
  % sigma: two lobes overlapping head-on
  \fill[omDef!35, opacity=0.8] (0,0) ellipse (0.9 and 0.42);
  \fill[omProp!35, opacity=0.8] (1.1,0) ellipse (0.9 and 0.42);
  \fill (0.1,0) circle (1.6pt); \fill (1.0,0) circle (1.6pt);
  \draw[dashed, black!50] (-1.2,0) -- (2.3,0);
  \node[below] at (0.55,-1.05) {$\sigma$：头对头，在轴线上};
  % pi: two p orbitals side by side
  \begin{scope}[xshift=5.2cm]
    \foreach \x/\c in {0/omDef, 1.2/omProp} {
      \fill[\c!35, opacity=0.85] (\x,0.45) ellipse (0.32 and 0.45);
      \fill[\c!15, opacity=0.85] (\x,-0.45) ellipse (0.32 and 0.45);
      \draw[\c!70!black] (\x,0.45) ellipse (0.32 and 0.45) (\x,-0.45) ellipse (0.32 and 0.45);
      \fill (\x,0) circle (1.6pt);
    }
    \draw[dashed, black!50] (-0.7,0) -- (1.9,0);
    \draw[omThm, thick, densely dotted] (0.25,0.75) -- (0.95,0.75) (0.25,-0.75) -- (0.95,-0.75);
    \node[below] at (0.6,-1.05) {$\pi$：侧面重叠，在轴线上下两侧};
  \end{scope}
\end{tikzpicture}

{\small 两种重叠方式。$\sigma$ 键允许绕其轴线旋转，
$\pi$ 键则不允许，这正是双键具有刚性的原因
（见\cref{ch:b1:stereochemistry-in-depth}）。}
\end{omfigure}

\section{价层电子对互斥模型}

\begin{definition}[价层电子对互斥模型、空间数]\label{def:b1:lewis-resonance-vsepr:vsepr}
在\emph{价层电子对互斥模型}\index{价层电子对互斥模型}中，
中心原子 A 周围的电子对彼此排斥，
并取使它们相距最远的排列方式。分子记为
\ce{AX_nE_m}，其中 $n$ 为与 A 成键的原子数
（多重键按一个计），$m$ 为 A 的\omterm{def:b1:lewis-resonance-vsepr:lewis}{孤电子对}数。
\emph{空间数}\index{空间数} $n + m$ 决定
电子对的排列方式：2 为直线形，3 为平面三角形，4 为四面体形，5
为三角双锥形，6 为八面体形。分子的形状
仅由其原子的排列决定。
\end{definition}

\begin{proposition}[VSEPR 几何构型]\label{prop:b1:lewis-resonance-vsepr:geometries}
\omterm{def:b1:lewis-resonance-vsepr:vsepr}{价层电子对互斥模型}可预测下表中的形状与理想键角；
实测键角与之相近。
\begin{center}
\small
\begin{tabular}{llllc}
\hline
类型 & 电子对排列 & 分子形状 & 例子 & 实测键角 \\
\hline
\ce{AX2} & 直线形 & 直线形 & \ce{CO2} & $180^\circ$ \\
\ce{AX3} & 平面三角形 & 平面三角形 & \ce{BF3} & $120^\circ$ \\
\ce{AX2E} & 平面三角形 & V 形 & \ce{SO2} & $119.5^\circ$ \\
\ce{AX4} & 四面体形 & 四面体形 & \ce{CH4} & $109.5^\circ$ \\
\ce{AX3E} & 四面体形 & 三角锥形 & \ce{NH3} & $106.7^\circ$ \\
\ce{AX2E2} & 四面体形 & V 形 & \ce{H2O} & $104.5^\circ$ \\
\ce{AX5} & 三角双锥形 & 三角双锥形 & \ce{PCl5} & $90^\circ$，$120^\circ$ \\
\ce{AX4E} & 三角双锥形 & 跷跷板形 & \ce{SF4} & $101.6^\circ$，$173.1^\circ$ \\
\ce{AX3E2} & 三角双锥形 & T 形 & \ce{ClF3} & $87.5^\circ$ \\
\ce{AX2E3} & 三角双锥形 & 直线形 & \ce{XeF2} & $180^\circ$ \\
\ce{AX6} & 八面体形 & 八面体形 & \ce{SF6} & $90^\circ$ \\
\ce{AX5E} & 八面体形 & 四方锥形 & \ce{BrF5} &——\\
\ce{AX4E2} & 八面体形 & 平面正方形 & \ce{XeF4} &——\\
\hline
\end{tabular}
\end{center}
% ledger: geo:CO2.aOCO, geo:BF3.aFBF, geo:SO2.aOSO, geo:CH4.aHCH,
% geo:NH3.aHNH, geo:H2O.aHOH, geo:SF6.aFSF, geo:SF4.aFSF2, geo:SF4.aFSF,
% geo:ClF3.aFClF, geo:XeF2.aFXeF
\end{proposition}
\admitted

\begin{proposition}[孤电子对使键角变小]\label{prop:b1:lewis-resonance-vsepr:lone-pair-angles}
\omterm{def:b1:lewis-resonance-vsepr:lewis}{孤电子对}仅由一个原子核束缚，其占据的空间比\omterm{def:b1:lewis-resonance-vsepr:lewis}{成键电子对}更大，
排斥也更强：排斥作用依次减弱的顺序为
\omterm{def:b1:lewis-resonance-vsepr:lewis}{孤电子对}/\omterm{def:b1:lewis-resonance-vsepr:lewis}{孤电子对} $>$ \omterm{def:b1:lewis-resonance-vsepr:lewis}{孤电子对}/\omterm{def:b1:lewis-resonance-vsepr:lewis}{成键电子对} $>$ \omterm{def:b1:lewis-resonance-vsepr:lewis}{成键电子对}/\omterm{def:b1:lewis-resonance-vsepr:lewis}{成键电子对}。
因此，随着\omterm{def:b1:lewis-resonance-vsepr:lewis}{孤电子对}增加，键角逐渐变小
（\ce{CH4} 中为 $109.5^\circ$，\ce{NH3} 中为 $106.7^\circ$，
\ce{H2O} 中为 $104.5^\circ$）；在三角双锥中，\omterm{def:b1:lewis-resonance-vsepr:lewis}{孤电子对}占据平伏位置，
在该位置上它们仅与两个相邻原子成 $90^\circ$ 夹角。
% ledger: geo:CH4.aHCH, geo:NH3.aHNH, geo:H2O.aHOH
\end{proposition}
\admitted

\begin{omfigure}
\setchemfig{atom sep=2.6em}
\renewcommand{\arraystretch}{1.2}
\begin{tabular}{ccccc}
\chemfig{O=C=O} &
\chemfig{B(-[:90]F)(-[:210]F)(-[:330]F)} &
\chemfig{C(-[:90]H)(-[:200]H)(<[:280]H)(<:[:335]H)} &
\chemfig{P(-[:90]Cl)(-[:270]Cl)(-[:180]Cl)(<[:320]Cl)(<:[:30]Cl)} &
\chemfig{S(-[:90]F)(-[:270]F)(-[:180]F)(-[:0]F)(<[:225]F)(<:[:45]F)} \\
\small \ce{AX2}, \ce{CO2} & \small \ce{AX3}, \ce{BF3} & \small \ce{AX4}, \ce{CH4} &
\small \ce{AX5}, \ce{PCl5} & \small \ce{AX6}, \ce{SF6} \\[16pt]
\chemfig{\charge{90=\:}{S}(=[:210]O)(=[:330]O)} &
\chemfig{\charge{90=\:}{N}(-[:200]H)(<[:280]H)(<:[:335]H)} &
\chemfig{\charge{60=\:,120=\:}{O}(-[:215]H)(-[:325]H)} &
\chemfig{\charge{0=\:}{S}(-[:90]F)(-[:270]F)(<[:210]F)(<:[:150]F)} &
\chemfig{Xe(-[:0]F)(-[:90]F)(-[:180]F)(-[:270]F)} \\
\small \ce{AX2E}, \ce{SO2} & \small \ce{AX3E}, \ce{NH3} & \small \ce{AX2E2}, \ce{H2O} &
\small \ce{AX4E}, \ce{SF4} & \small \ce{AX4E2}, \ce{XeF4} \\
\end{tabular}

\medskip
{\small VSEPR 预测的形状。楔形键指向读者，
虚楔形键背离读者；普通键位于纸面内。
中心原子的孤对电子用成对的点表示；\ce{XeF4} 为正面图，
其两对孤对电子分别指向读者与背离读者。}
\end{omfigure}

\begin{method}[预测形状]\label{met:b1:lewis-resonance-vsepr:vsepr}
预测中心原子周围的形状，步骤如下：
\begin{enumerate}
  \item 画出\omterm{def:b1:lewis-resonance-vsepr:lewis}{路易斯结构}（见\cref{met:b1:lewis-resonance-vsepr:lewis}）；
  \item 数出与中心原子成键的原子数（$n$）及其孤对
        电子对数（$m$），记为 \ce{AX_nE_m}；
  \item 由 $n + m$ 得出排布，由
        $n$ 个原子的位置得出形状；
  \item 修正理想键角：孤对电子与多重键占据的空间
        比单键更大。
\end{enumerate}
\end{method}

\begin{definition}[杂化]\label{def:b1:lewis-resonance-vsepr:hybridisation}
原子的\emph{杂化}\index{杂化}是对其化学键的一种
描述，即用沿 VSEPR 所预测方向伸展的轨道表示成键：
\omterm{def:b1:lewis-resonance-vsepr:vsepr}{空间数}为 4 的原子称为 sp$^3$ 杂化（四个等价轨道，
夹角 $109.5^\circ$）；\omterm{def:b1:lewis-resonance-vsepr:vsepr}{空间数}为 3 的为 sp$^2$ 杂化（三个轨道
共面、夹角 $120^\circ$，另留一个 p 轨道
与之垂直）；\omterm{def:b1:lewis-resonance-vsepr:vsepr}{空间数}为 2 的为 sp 杂化（两个轨道，夹角 $180^\circ$，
另留两个 p 轨道）。剩余的 p 轨道形成 $\pi$ 键。
\end{definition}

\begin{example}[有机化学中的碳原子]\label{ex:b1:lewis-resonance-vsepr:carbon}
乙烷中每个碳原子均为 sp$^3$，呈四面体形；乙烯中为 sp$^2$，六个
原子共面，键角接近 $120^\circ$（实测 \ce{H-C-H}
为 $117.6^\circ$）；乙炔为 sp，四个原子共线。键
长由 153.6 缩短至 133.9，再缩至 \qty{120.3}{pm}，
这是 $\pi$ 键增加的结果。
% ledger: geo:C2H4.aHCH, geo:C2H6.rCC, geo:C2H4.rCC, geo:C2H2.rCC
\end{example}

\begin{remark}[描述而非原因]\label{rem:b1:lewis-resonance-vsepr:hybrid-limit}
\omterm{def:b1:lewis-resonance-vsepr:hybridisation}{杂化}描述的是已知的几何构型，它并不解释
其成因；它对电子能量也无能为力，这需要由分子轨道理论
来解释，见本科二年级分册。\omterm{def:b1:lewis-resonance-vsepr:hybridisation}{杂化}仍是有机化学的日常用语，
“sp$^2$ 碳”指平面形的碳，其 p 轨道
可用于形成 $\pi$ 键。
\end{remark}

\section{极性分子：偶极矩}

\begin{definition}[键偶极，偶极矩]\label{def:b1:lewis-resonance-vsepr:dipole}
在\omterm{def:b1:periodicity:electronegativity}{电负性}不同的两原子所成的键中，电子
偏向\omterm{def:b1:periodicity:electronegativity}{电负性}较大的原子，该原子带部分电荷，
$-\delta e$，另一原子带 $+\delta e$（$0 < \delta < 1$）。
\emph{键偶极}\index{键偶极}是矢量 $\vec\mu = \delta e\,
\vec d$，模为 $\delta e d$，方向由负
部分电荷指向正部分电荷，$\vec d$ 为两者之间的矢量。
分子的\emph{偶极矩}\index{偶极矩}是其各键偶极
的矢量和（计入孤对电子的贡献）；偶极矩可用
库仑米或德拜表示，$\qty{1}{\debye} =
\qty{3.336e-30}{C.m}$。偶极矩不为零的分子是
\emph{极性分子}\index{极性分子}。
% ledger: const:c (1 D = 1e-21/c C.m)
\end{definition}

\begin{remark}[方向约定]\label{rem:b1:lewis-resonance-vsepr:convention}
物理学中偶极矢量由负电荷指向正电荷，
即本书所采用的约定。较旧的化学书中则画
一端带叉尾的箭头、指向负端；两者仅
方向不同。
\end{remark}

\begin{proposition}[对称分子的偶极矩]\label{prop:b1:lewis-resonance-vsepr:dipole-sum}
中心原子 A 上无孤对电子的分子 \ce{AX_n}（直线形 \ce{AX2}、平面三角形
\ce{AX3}、四面体形 \ce{AX4}、\dots）的\omterm{def:b1:lewis-resonance-vsepr:dipole}{偶极矩}为零，无论
其键的极性如何。键角为 $\theta$ 的角形 \ce{AX2} 分子，
若其\omterm{def:b1:lewis-resonance-vsepr:dipole}{键偶极}为 $\mu_b$，则\omterm{def:b1:lewis-resonance-vsepr:dipole}{偶极矩}为
\[
  \mu = 2\mu_b \cos\frac{\theta}{2} .
\]
\end{proposition}

\begin{proof}
对称构型中，各键矢量的模相等，且其
方向关于 A 对称分布，因此矢量和为零
（直线形 \ce{AX2} 中二者方向相反；\ce{AX3} 中夹角为 $120^\circ$，
三个等模矢量之和为零；\ce{AX4} 中，和式
在四面体的旋转下保持不变，故其为零）。对于
角形分子，每个键矢量与角平分线的夹角为 $\theta/2$，
垂直于角平分线的分量相互抵消，
沿角平分线的分量相加，得 $2\mu_b\cos(\theta/2)$。
\end{proof}

\begin{omfigure}
\begin{tikzpicture}[font=\small, >={Stealth}]
  % water
  \node[aO] (o) at (0,0) {};
  \node[aH] (h1) at (-0.95,-0.73) {};
  \node[aH] (h2) at (0.95,-0.73) {};
  \begin{scope}[on background layer]
    \draw[bond] (o.center) -- (h1.center); \draw[bond] (o.center) -- (h2.center);
  \end{scope}
  \draw[->, omProp, thick] (-0.25,0.15) -- (-0.95,-0.38);
  \draw[->, omProp, thick] (0.25,0.15) -- (0.95,-0.38);
  \draw[->, omThm, very thick] (0,-0.25) -- (0,-1.45);
  \node[omThm, right] at (0.05,-1.25) {$\vec\mu$};
  \node at (0,0.55) {$\delta^-$};
  \node at (-1.25,-1.05) {$\delta^+$}; \node at (1.25,-1.05) {$\delta^+$};
  \node[below] at (0,-1.7) {\ce{H2O}：$\mu = \qty{1.86}{\debye}$};
  % carbon dioxide
  \begin{scope}[xshift=5cm]
    \node[aO] (a) at (-1.1,0) {}; \node[aC] (c) at (0,0) {}; \node[aO] (b) at (1.1,0) {};
    \begin{scope}[on background layer]
      \foreach \d in {-1.6,1.6} {
        \draw[bond, line width=1.1pt] ([yshift=\d pt]a.center) -- ([yshift=\d pt]c.center);
        \draw[bond, line width=1.1pt] ([yshift=\d pt]c.center) -- ([yshift=\d pt]b.center);}
    \end{scope}
    \draw[->, omProp, thick] (-0.75,0.45) -- (-0.2,0.45);
    \draw[->, omProp, thick] (0.75,0.45) -- (0.2,0.45);
    \node[below] at (0,-1.7) {\ce{CO2}：键偶极相互抵消，$\mu = 0$};
  \end{scope}
\end{tikzpicture}

{\small \omterm{def:b1:lewis-resonance-vsepr:dipole}{键偶极}（橙色，从部分负
电荷指向部分正电荷）及其矢量和（红色）。在折线形的水中二者相加；在直线形的
二氧化碳中二者相互抵消。}
% ledger: dip:H2O
\end{omfigure}

\begin{example}[极性分子与非极性分子]\label{ex:b1:lewis-resonance-vsepr:dipoles}
\ce{CO2}、\ce{BF3}、\ce{CH4} 和 \ce{CCl4} 的\omterm{def:b1:lewis-resonance-vsepr:dipole}{偶极矩}均为零。
\ce{H2O}（\qty{1.86}{\debye}）、\ce{NH3}（\qty{1.48}{\debye}）和
\ce{SO2}（\qty{1.63}{\debye}）均为\omterm{def:b1:lewis-resonance-vsepr:dipole}{极性分子}。沿 \ce{CH3Cl}、\ce{CH2Cl2}、
\ce{CHCl3} 一列，\omterm{def:b1:lewis-resonance-vsepr:dipole}{偶极矩}依次降低，分别为 1.87、1.62、\qty{1.04}{\debye}，
至 \ce{CCl4} 中降为零：\ce{C-Cl} \omterm{def:b1:lewis-resonance-vsepr:dipole}{键偶极}越来越相互抵消。
在卤化氢中，\omterm{def:b1:lewis-resonance-vsepr:dipole}{偶极矩}随\omterm{def:b1:periodicity:electronegativity}{电负性}差而变化：HF 1.83、
HCl 1.09、HBr 0.83、HI \qty{0.45}{\debye}。
% ledger: dip:H2O, dip:NH3, dip:SO2, dip:CH3Cl, dip:CH2Cl2, dip:CHCl3,
% dip:CCl4, dip:HF, dip:HCl, dip:HBr, dip:HI
\end{example}

\section{练习}

\begin{exercise}[$\star$]\label{exo:b1:lewis-resonance-vsepr:1}
画出下列分子的\omterm{def:b1:lewis-resonance-vsepr:lewis}{路易斯结构}：\ce{H2O}、\ce{NH3}、\ce{CO2}、\ce{HCN}、
\ce{CH2O}（甲醛）和 \ce{NH4+}，并给出\omterm{def:b1:lewis-resonance-vsepr:formal-charge}{形式电荷}。
\end{exercise}

\begin{exercise}[$\star$]\label{exo:b1:lewis-resonance-vsepr:2}
给出下列分子的 \ce{AX_nE_m} 类型、形状和近似键角：
\ce{H2S}、\ce{PCl3}、\ce{BF3}、\ce{SiH4} 和 \ce{CO3^{2-}}（每式中的
中心原子列在最前）。
\end{exercise}

\begin{exercise}[$\star$]\label{exo:b1:lewis-resonance-vsepr:3}
下列分子中哪些是\omterm{def:b1:lewis-resonance-vsepr:dipole}{极性分子}：\ce{CCl4}、\ce{CH2Cl2}、\ce{BF3}、
\ce{NH3}、\ce{CO2}、\ce{SO2}？请根据分子形状说明理由。
\end{exercise}

\begin{exercise}[$\star$]\label{exo:b1:lewis-resonance-vsepr:4}
指出下列反应中的\omterm{def:b1:lewis-resonance-vsepr:lewis-acid}{路易斯酸}与\omterm{def:b1:lewis-resonance-vsepr:lewis-acid}{路易斯碱}：\ce{H+ + H2O -> H3O+} 及
\ce{Cu^{2+} + 4NH3 -> [Cu(NH3)4]^{2+}}。哪些键是\omterm{def:b1:lewis-resonance-vsepr:lewis-acid}{配位键}？
\end{exercise}

\begin{exercise}[$\star\star$]\label{exo:b1:lewis-resonance-vsepr:5}
写出乙酸根离子 \ce{CH3COO-} 的两种\omterm{def:b1:lewis-resonance-vsepr:resonance}{共振结构}。它们
对两个 \ce{C-O} 键长有何预测？将其与甲酸
\ce{HCOOH} 作比较，后者两个 \ce{C-O} 键长分别为 120.2 和
\qty{134.3}{pm}。
% ledger: geo:H2CO2.rCO, geo:H2CO2.rCO2
\end{exercise}

\begin{exercise}[$\star\star$]\label{exo:b1:lewis-resonance-vsepr:6}
画出硫酸根离子 \ce{SO4^{2-}} 的\omterm{def:b1:lewis-resonance-vsepr:lewis}{路易斯结构}，其中含四个
\ce{S-O} 单键，计算\omterm{def:b1:lewis-resonance-vsepr:formal-charge}{形式电荷}，并预测
其形状。可写出多少种含两个 \ce{S=O} 双键、
且硫原子上不带电荷的\omterm{def:b1:lewis-resonance-vsepr:resonance}{共振结构}？
\end{exercise}

\begin{exercise}[$\star\star$]\label{exo:b1:lewis-resonance-vsepr:7}
用\omterm{def:b1:lewis-resonance-vsepr:vsepr}{价层电子对互斥模型}预测 \ce{SF4}、\ce{ClF3}、\ce{XeF2} 和
\ce{XeF4} 的形状。试说明三角双锥中\omterm{def:b1:lewis-resonance-vsepr:lewis}{孤电子对}为何位于
平伏位所在的平面内。
\end{exercise}

\begin{exercise}[$\star\star$]\label{exo:b1:lewis-resonance-vsepr:8}
键角 \ce{H-X-H} 在 \ce{H2O} 中为 $104.5^\circ$，在
\ce{H2S} 中为 $92.1^\circ$；在 \ce{NH3} 中为 $106.7^\circ$，在 \ce{PH3} 中为 $93.3^\circ$。
试从中心原子的大小与\omterm{def:b1:periodicity:electronegativity}{电负性}出发，
给出这些键角差异的解释。
% ledger: geo:H2O.aHOH, geo:H2S.aHSH, geo:NH3.aHNH, geo:PH3.aHPH
\end{exercise}

\begin{exercise}[$\star\star$]\label{exo:b1:lewis-resonance-vsepr:9}
给出下列分子中各碳原子及氮原子的\omterm{def:b1:lewis-resonance-vsepr:hybridisation}{杂化}类型，
即 \ce{CH3-C#N}（乙腈）和 \ce{CH2=CH-CH3}，并给出每个分子中的
$\sigma$ 键和 $\pi$ 键的数目。
\end{exercise}

\begin{exercise}[$\star\star\star$]\label{exo:b1:lewis-resonance-vsepr:10}
\ce{NF3} 的\omterm{def:b1:lewis-resonance-vsepr:dipole}{偶极矩}仅为 \qty{0.24}{\debye}，而
\ce{NH3} 为 \qty{1.48}{\debye}，尽管 \ce{N-F} 键的极性比
\ce{N-H} 键更强。利用\omterm{def:b1:lewis-resonance-vsepr:dipole}{键偶极}与氮原子上的\omterm{def:b1:lewis-resonance-vsepr:lewis}{孤电子对}，
解释二者的差异。
% ledger: dip:NF3, dip:NH3
\end{exercise}

\begin{exercise}[$\star\star\star$]\label{exo:b1:lewis-resonance-vsepr:11}
\ce{SO2} 分子的键角为 $119.5^\circ$，\omterm{def:b1:lewis-resonance-vsepr:dipole}{偶极矩}为
\qty{1.63}{\debye}。计算一根 \ce{S-O} 键的\omterm{def:b1:lewis-resonance-vsepr:dipole}{偶极矩}；已知键长
为 \qty{143.2}{pm}，推出
每个氧原子上的部分电荷 $\delta$。
% ledger: geo:SO2.aOSO, geo:SO2.rSO, dip:SO2
\end{exercise}

\begin{exercise}[$\star\star\star$]\label{exo:b1:lewis-resonance-vsepr:12}
画出一氧化二氮 \ce{N2O}（直线形，\ce{N-N-O}）的最重要\omterm{def:b1:lewis-resonance-vsepr:lewis}{路易斯结构}，
并计算\omterm{def:b1:lewis-resonance-vsepr:formal-charge}{形式电荷}。实测键长 \ce{N-N} 为 \qty{112.8}{pm}，\ce{N-O} 为
\qty{118.4}{pm}；在 \ce{N2} 中三键的键长为 \qty{109.8}{pm}。
哪一种结构的
贡献更大？
% ledger: geo:N2O.rNN, geo:N2O.rNO, re:N2
\end{exercise}

\section{问题：雷击中的分子}

\begin{problem}[{周末问题——闪电加热空气时生成的氮、氧物种：
路易斯结构、共振、形状，以及水的
部分电荷}]
\label{pb:b1:lewis-resonance-vsepr:1}
一道闪电将空气加热至数千摄氏度：\ce{N2} 与 \ce{O2}
生成 \ce{NO}，随后生成 \ce{NO2}、\ce{N2O}、臭氧，最后在
雨水中形成硝酸。数据：$\qty{1}{\debye} = \qty{3.336e-30}{C.m}$，$e =
\qty{1.602e-19}{C}$。需要时给出实测值。

\medskip
\noindent\textbf{第一部分——\omterm{def:b1:lewis-resonance-vsepr:lewis}{路易斯结构}。}
\begin{enumerate}
  \item 数出 \ce{N2}、\ce{NO}、\ce{NO2}、
        \ce{N2O}、\ce{O3} 和 \ce{HNO3} 的\omterm{def:b1:quantum-numbers:configuration}{价电子}数。
  \item 画出 \ce{N2} 的\omterm{def:b1:lewis-resonance-vsepr:lewis}{路易斯结构}。
  \item 画出 \ce{NO} 的\omterm{def:b1:lewis-resonance-vsepr:lewis}{路易斯结构}。为什么
        两个原子不能同时满足\omterm{def:b1:lewis-resonance-vsepr:lewis}{八隅规则}？
  \item 画出 \ce{N2O} 的两种\omterm{def:b1:lewis-resonance-vsepr:lewis}{路易斯结构}，分别为 \ce{N=N=O} 与
        \ce{N#N-O}，并给出\omterm{def:b1:lewis-resonance-vsepr:formal-charge}{形式电荷}。
  \item 画出臭氧的\omterm{def:b1:lewis-resonance-vsepr:lewis}{路易斯结构}，并给出\omterm{def:b1:lewis-resonance-vsepr:formal-charge}{形式电荷}。
  \item 画出 \ce{HNO3} 的\omterm{def:b1:lewis-resonance-vsepr:lewis}{路易斯结构}，并给出\omterm{def:b1:lewis-resonance-vsepr:formal-charge}{形式电荷}。
  \item 六种物种中，哪些含有\omterm{def:b1:quantum-numbers:unpaired}{未成对电子}？
\end{enumerate}

\medskip
\noindent\textbf{第二部分——共振与键长。}
\begin{enumerate}[resume]
  \item 写出臭氧的两种\omterm{def:b1:lewis-resonance-vsepr:resonance}{共振结构}。它们预测的
        键级是多少？
  \item 臭氧中 \ce{O-O} 键长为 \qty{127.8}{pm}，
        \ce{O2} 中为 \qty{120.8}{pm}，\ce{H2O2} 中为 \qty{147.5}{pm}。
        该预测是否得到证实？
  \item 写出 \ce{NO3-} 的三种\omterm{def:b1:lewis-resonance-vsepr:resonance}{共振结构}，并给出
        每个 \ce{N-O} 键的键级。
  \item \ce{HNO3} 中三个 \ce{N-O} 键的键长分别为 140.6、121.1 和
        \qty{119.9}{pm}。请将它们分别指认到各键上，并用共振加以解释。
  \item 解释 \ce{O-N-O} 键角（\ce{NO2} 中为 $134^\circ$）为何
        大于 $120^\circ$。
  \item 第 4 题的两种结构中，哪一种把负\omterm{def:b1:lewis-resonance-vsepr:formal-charge}{形式电荷}放在
        \omterm{def:b1:periodicity:electronegativity}{电负性}更大的原子上？
\end{enumerate}

\medskip
\noindent\textbf{第三部分——形状与极性。}
\begin{enumerate}[resume]
  \item 给出中心原子的类型 \ce{AX_nE_m}，分别针对 \ce{O3}、\ce{N2O}、\ce{NO3-}，
        以及 \ce{HNO3} 中的氮原子，并说明它们的形状。
  \item 臭氧的键角为 $116.8^\circ$，试解释它为何小于
        $120^\circ$。
  \item 臭氧的\omterm{def:b1:lewis-resonance-vsepr:dipole}{偶极矩}为 \qty{0.53}{\debye}，\ce{N2O} 为
        \qty{0.16}{\debye}，\ce{N2} 为零。请解释。
  \item 为什么 \ce{CO2} 是非极性的，而 \ce{SO2} 是极性的？
  \item 预测 \ce{NH3} 的键角，并与 $109.5^\circ$ 比较，
        再与实测值 $106.7^\circ$ 进行对照。
  \item 对水（$104.5^\circ$）提出同样的问题。
\end{enumerate}
% ledger: geo:O3.rOO, re:O2, geo:H2O2.rOO, geo:HNO3.rNO, geo:HNO3.rNO2,
% geo:HNO3.rNO3, geo:NO2.aONO, geo:O3.aOOO, dip:O3, dip:N2O

\medskip
\noindent\textbf{第四部分——水的部分电荷。}
水的\omterm{def:b1:lewis-resonance-vsepr:dipole}{偶极矩}为 \qty{1.857}{\debye}，其键角为
$104.48^\circ$，\ce{O-H} 键长为 \qty{95.8}{pm}。
% ledger: dip:H2O, geo:H2O.aHOH, geo:H2O.rOH
\begin{enumerate}[resume]
  \item 把水的\omterm{def:b1:lewis-resonance-vsepr:dipole}{偶极矩}表示为\omterm{def:b1:lewis-resonance-vsepr:dipole}{键偶极}
        $\mu_{OH}$ 与键角的函数。
  \item 以德拜为单位，计算 $\mu_{OH}$。
  \item 将其换算为库仑米。
  \item 若\omterm{def:b1:lewis-resonance-vsepr:dipole}{键偶极}由位于两个原子核上的两个满电荷 $\pm e$ 构成，
        它将是多少？
  \item 推导 \ce{O-H} 键的离子性分数 $\delta$，
        即实测\omterm{def:b1:lewis-resonance-vsepr:dipole}{键偶极}与满电荷对应的
        \omterm{def:b1:lewis-resonance-vsepr:dipole}{键偶极}之比。
\end{enumerate}
\end{problem}
